\section{Global assembly and coefficient}

At $r=r_*$, \cref{eq:endpointreduction,eq:crossing} give
$M(r_*)=q$.  Therefore \cref{eq:surfaceuppergeneral} becomes
\begin{equation}\label{eq:surfaceupper}
 \cH^2(\partial U_{r_*})
 \le4\pi r_*^2qD.
\end{equation}
Combining \cref{eq:surfacelower,eq:surfaceupper} yields
\[
 D\ge
 \frac{(18/\pi^2)^{1/3}}{r_*^2q}n^{2/3}.
\]
Using \cref{eq:localidentity},
\begin{equation}\label{eq:Dexact}
 \boxed{
 D\ge
 \left(\frac{18}{\pi^2}\right)^{1/3}
 \left(103-\frac{117\sqrt3}{2}\right)n^{2/3}.}
\end{equation}
For a contact maximizer, $|E|=f_3(n)$, proving \cref{eq:mainexact}.

To prove the clean decimal, use the rigorous bounds
\[
 \sqrt3<\frac{1732050808}{10^9},
 \qquad
 \pi<\frac{3141593}{10^6}.
\]
The first gives
\[
 103-\frac{117\sqrt3}{2}
 >\frac{418756933}{250000000}.
\]
For $c_0=4093/2000=2.0465$, exact rational arithmetic gives
\begin{align*}
&18\left(\frac{418756933}{250000000}\right)^3
-c_0^3\left(\frac{3141593}{10^6}\right)^2\\
&\qquad=
\frac{83319151211028177098003}
{125000000000000000000000000}>0.
\end{align*}
All factors are positive, hence $c_0<c_K$.  This proves \cref{eq:clean} and the
strict theorem.

\section{Verification and scope}

The finite theorem uses the following named published inputs:
\begin{enumerate}[label=(\roman*),leftmargin=2.2em]
\item the three-dimensional kissing number;
\item the Kepler density $\delta_3=\pi/\sqrt{18}$;
\item Euclidean isoperimetry for bounded finite-perimeter sets;
\item spherical-neighborhood isoperimetry for measurable subsets of $\Sph^2$.
\end{enumerate}
The finite outer-parallel inequality itself is proved in
\cref{prop:finitekepler} from the explicit FCC lattice, translation averaging on
a compact flat torus, and input (ii).  In particular, the theorem does not depend
on Bezdek--L\'angi equation (1.2) or on any product rule for separate limsups.

The accompanying exact verifier checks the optimized-radius identities in
$\mathbb Q(\sqrt3)$, rigorous rational bounds for $\sqrt3$ and $\pi$, and the
positive cubed margin proving $c_K>2.0465$.  The public Lean development checks the
finite contact-model relations, the optimized local algebra, the degree-sum and
surface-assembly implications, and the relational conclusion for $f_3(n)$ from
explicit geometric-input structures.  It does not yet reconstruct the flat-torus
completion proof or foundational versions of (i)--(iv); ordinary theorem
parameters remain visible in the load-bearing signatures.

No Barlow, crystallization, $\Gamma$-convergence, interface, or Wulff-selection
result enters the proof.  The existence and exact value of
\[
 \lim_{n\to\infty}\frac{6n-f_3(n)}{n^{2/3}}
\]
remain open.

\begin{thebibliography}{99}

\bibitem{BaernsteinTaylor1976}
A.~Baernstein II and B.~A.~Taylor,
\emph{Spherical rearrangements, subharmonic functions, and $*$-functions in
$n$-space}, Duke Math. J. \textbf{43} (1976), 245--268.

\bibitem{BarnesOzgurWu2018}
L.~P.~Barnes, A.~\"Ozg\"ur, and X.~Wu,
\emph{An isoperimetric result on high-dimensional spheres},
arXiv:1811.10533 (2018), Proposition~1.

\bibitem{BezdekKhan2018}
K.~Bezdek and M.~A.~Khan,
\emph{Contact numbers for sphere packings},
in New Trends in Intuitive Geometry,
Bolyai Soc. Math. Stud. \textbf{27}, Springer, 2018, 25--47.

\bibitem{BezdekLangi2025}
K.~Bezdek and Z.~L\'angi,
\emph{Density bounds for unit ball packings relative to their outer parallel
domains}, Pure Appl. Funct. Anal. \textbf{10} (2025), 1193--1205;
arXiv:2401.00645.

\bibitem{BezdekReid2013}
K.~Bezdek and S.~Reid,
\emph{Contact graphs of unit sphere packings revisited},
J. Geom. \textbf{104} (2013), 57--83.

\bibitem{Erdos1975}
P.~Erd\H{o}s,
\emph{On some problems of elementary and combinatorial geometry},
Ann. Mat. Pura Appl. (4) \textbf{103} (1975), 99--108.

\bibitem{Hales2005}
T.~C.~Hales,
\emph{A proof of the Kepler conjecture},
Ann. of Math. (2) \textbf{162} (2005), 1065--1185.

\bibitem{Maggi2012}
F.~Maggi,
\emph{Sets of Finite Perimeter and Geometric Variational Problems},
Cambridge Studies in Advanced Mathematics 135, Cambridge University Press, 2012.

\bibitem{SchuetteWaerden1953}
K.~Sch\"utte and B.~L.~van der Waerden,
\emph{Das Problem der dreizehn Kugeln},
Math. Ann. \textbf{125} (1952/53), 325--334.

\end{thebibliography}
